\section{Fourier Analysis: From Time to Frequency}
\label{sec:fourier}

The accelerometer hands us a wiggling line: the magnitude \(a[n]\) sampled
\SI{64}{\hertz}, one number every \(\approx\SI{15.6}{\milli\second}\). Looking at
that line, a calm stroll and a freeze episode are both ``a wobble around \(1\,g\)''
--- the eye cannot reliably separate them. The whole trick of Cadence is to stop
asking \emph{how big} the wobble is and start asking \emph{how fast} it wobbles.
Walking paces the leg at roughly \(1\text{--}2\) steps per second; a freeze sets the
foot trembling in place at a much faster \(3\text{--}8\) times per second. Those two
\emph{rhythms} live in different places on the frequency axis even when their
amplitudes overlap. Fourier analysis is the lens that turns a tangle of time-domain
wiggles into a clean picture of which rhythms are present and how strong each one is.

\begin{intuition}
A chord on a piano is a single complicated pressure wave hitting your ear, yet you
hear three distinct notes. Your ear is doing Fourier analysis: it decomposes one
messy waveform into the pure tones that add up to it. Cadence does the same to the
ankle signal --- it listens for the ``note'' of a freeze (a \(3\text{--}8\,\)Hz
buzz) hiding underneath the ``note'' of walking (a \(1\text{--}2\,\)Hz sway).
\end{intuition}

\subsection{Any signal is a sum of sinusoids}

The founding idea, due to Joseph Fourier, is that \emph{any} reasonable signal --- no
matter how jagged --- can be rebuilt by adding together pure sine and cosine waves of
different frequencies, each with its own amplitude and phase. A slow sine captures the
broad sway; faster sines fill in the finer trembles. Add enough of them and you
reproduce the original wiggle exactly. The two questions ``what does the signal look
like in time?'' and ``which sinusoids is it made of?'' are two descriptions of the
\emph{same} object --- the second is just the more useful one for telling walking
apart from freezing.

\begin{figure}[H]
\centering
\resizebox{0.96\textwidth}{!}{%
\begin{tikzpicture}
% ----- LEFT: time-domain signal = sum of two sinusoids -----
\begin{axis}[
  name=time,
  width=8.8cm,height=5.4cm,
  xlabel={time \(t\) (s)}, ylabel={\(a(t)\) (mg, DC removed)},
  axis lines=left, xmin=0,xmax=4, ymin=-2.6,ymax=2.6,
  xtick={0,1,2,3,4}, ytick=\empty,
  every axis label/.append style={font=\small},
  tick label style={font=\footnotesize},
  title={\small\bfseries\color{ink} Time domain: one tangled wiggle},
]
% slow walking component
\addplot[steel,thick,dashed,domain=0:4,samples=300]{1.5*sin(deg(2*pi*1.5*x))};
% fast freeze component
\addplot[freeze,thick,dashdotted,domain=0:4,samples=400]{0.8*sin(deg(2*pi*5*x))};
% their sum
\addplot[navy,very thick,domain=0:4,samples=600]
  {1.5*sin(deg(2*pi*1.5*x)) + 0.8*sin(deg(2*pi*5*x))};
\end{axis}
\node[steel,font=\scriptsize] at (time.north west) [xshift=18mm,yshift=-3mm] {slow sway};
\node[freeze,font=\scriptsize] at (time.north west) [xshift=18mm,yshift=-8mm] {fast tremble};
% ----- arrow -----
\node[right=2mm of time.east,font=\Large\bfseries,navy] (ft) {\(\xrightarrow{\;\text{DFT}\;}\)};
% ----- RIGHT: frequency-domain spectrum, two peaks -----
\begin{axis}[
  at={($(time.east)+(2.7cm,0)$)}, anchor=west,
  width=8.8cm,height=5.4cm,
  xlabel={frequency \(f\) (Hz)}, ylabel={power \(|X(f)|^2\)},
  axis lines=left, xmin=0,xmax=10, ymin=0,ymax=1.15,
  xtick={0,1.5,3,5,8,10}, ytick=\empty,
  every axis label/.append style={font=\small},
  tick label style={font=\footnotesize},
  title={\small\bfseries\color{ink} Frequency domain: two clean peaks},
]
% locomotor band shading
\addplot[draw=none,fill=palenavy] coordinates {(0.5,0)(0.5,1.15)(3,1.15)(3,0)}\closedcycle;
% freeze band shading
\addplot[draw=none,fill=palefreeze] coordinates {(3,0)(3,1.15)(8,1.15)(8,0)}\closedcycle;
% walking peak at 1.5 Hz
\addplot[steel,very thick,domain=0:10,samples=400]{1.0*exp(-((x-1.5)^2)/0.05)};
% freeze peak at 5 Hz
\addplot[freeze,very thick,domain=0:10,samples=400]{0.45*exp(-((x-5)^2)/0.05)};
\node[steel,font=\scriptsize,align=center] at (axis cs:1.5,1.06) {walking\\\(1.5\,\)Hz};
\node[freeze,font=\scriptsize,align=center] at (axis cs:5,0.62) {freeze\\\(5\,\)Hz};
\node[navy,font=\scriptsize] at (axis cs:1.75,0.16) {loco band};
\node[freeze!70!black,font=\scriptsize] at (axis cs:5.5,0.16) {freeze band};
\end{axis}
\end{tikzpicture}}
\caption{The same signal, two views. \textbf{Left:} a slow \(1.5\,\)Hz sway
(steel, dashed) plus a faster \(5\,\)Hz tremble (red, dash-dotted) add to one tangled
trace (navy) --- hard to read by eye. \textbf{Right:} the Discrete Fourier Transform
unfolds that trace into a spectrum where each rhythm becomes a sharp peak at its own
frequency. Walking energy lands in the locomotor band (\(0.5\text{--}3\,\)Hz, navy
shading); freeze energy lands in the freeze band (\(3\text{--}8\,\)Hz, red shading).
Separating the two by \emph{rhythm} is what makes them distinguishable.}
\label{fig:time-to-freq}
\end{figure}

\subsection{The Discrete Fourier Transform and its bins}

We never have a continuous signal; we have \(N\) evenly spaced samples. The
\textbf{Discrete Fourier Transform} (DFT) is the finite-data version of Fourier's
idea. It takes a block of \(N\) samples and returns \(N\) complex numbers
\(X[k]\), each measuring how much of the \(k\)-th sinusoid the block contains.

\begin{mathsbox}
For a window of \(N\) samples \(a[0],\dots,a[N-1]\), the DFT coefficients are
\begin{equation}
X[k] \;=\; \sum_{n=0}^{N-1} a[n]\, e^{-\mathrm{j}\,2\pi kn/N},
\qquad k = 0,1,\dots,N-1 .
\label{eq:dft}
\end{equation}
The term \(e^{-\mathrm{j}2\pi kn/N}\) is a pure spinning sinusoid at frequency index
\(k\); the sum measures how strongly the data ``rings'' at that frequency.
Coefficient \(X[k]\) maps to a physical frequency
\begin{equation}
f_k \;=\; k\cdot\frac{f_s}{N}\,, \qquad k=0,1,\dots,\tfrac{N}{2}.
\label{eq:fk}
\end{equation}
The Fast Fourier Transform (FFT) is simply a clever \(\mathcal{O}(N\log N)\)
algorithm that computes the \emph{same} \(X[k]\) far faster than the naive
\(\mathcal{O}(N^2)\) sum --- the answer is identical.
\end{mathsbox}

The spectrum is not continuous: it is sampled at discrete frequencies \(f_k\) called
\textbf{bins}. The spacing between neighbouring bins is the single most important
number in this section:
\begin{equation}
\Delta f \;=\; \frac{f_s}{N} \;=\; \frac{64\ \text{Hz}}{256} \;=\; 0.25\ \text{Hz}.
\label{eq:binspacing}
\end{equation}
With \(f_s=\SI{64}{\hertz}\) and a window of \(N=256\) samples (\(=\SI{4.0}{\second}\)),
Cadence resolves the spectrum in steps of \(\SI{0.25}{\hertz}\) from \(0\) up to the
Nyquist limit \(f_s/2=\SI{32}{\hertz}\). That Nyquist ceiling comfortably covers every
band we care about (the highest is the \(8\,\)Hz top of the freeze band), so nothing of
interest is aliased.

\begin{figure}[H]
\centering
\begin{tikzpicture}
\begin{axis}[
  width=14.2cm,height=4.7cm,
  ybar interval=0.85,
  xlabel={frequency \(f\) (Hz) --- one bar per DFT bin, \(\Delta f=0.25\,\)Hz},
  ylabel={band power},
  axis lines=left, xmin=0,xmax=9, ymin=0,ymax=1.28,
  xtick={0,0.5,1,1.5,2,2.5,3,4,5,6,7,8},
  xticklabel style={font=\scriptsize}, ytick=\empty,
  every axis label/.append style={font=\small},
]
% draw bins as a coarse histogram approximating two band humps
\addplot[draw=steel,fill=palenavy,thick] coordinates {
 (0.5,0.10)(0.75,0.18)(1.0,0.34)(1.25,0.62)(1.5,0.92)(1.75,0.74)(2.0,0.46)
 (2.25,0.26)(2.5,0.15)(2.75,0.09)};
\addplot[draw=freeze,fill=palefreeze,thick] coordinates {
 (3.0,0.06)(3.25,0.10)(3.5,0.16)(3.75,0.22)(4.0,0.30)(4.25,0.40)(4.5,0.52)
 (4.75,0.58)(5.0,0.56)(5.25,0.44)(5.5,0.30)(5.75,0.20)(6.0,0.13)(6.25,0.09)
 (6.5,0.07)(6.75,0.06)(7.0,0.05)(7.25,0.045)(7.5,0.04)(7.75,0.04)};
% band boundary markers
\draw[densely dashed,navy] (axis cs:0.5,0) -- (axis cs:0.5,1.0);
\draw[densely dashed,navy] (axis cs:3,0) -- (axis cs:3,1.0);
\draw[densely dashed,freeze] (axis cs:8,0) -- (axis cs:8,1.0);
\node[navy,font=\scriptsize] at (axis cs:1.75,1.17) {locomotor band \(0.5\text{--}3\,\)Hz};
\node[freeze!75!black,font=\scriptsize] at (axis cs:5.6,1.17) {freeze band \(3\text{--}8\,\)Hz};
\end{axis}
\end{tikzpicture}
\caption{Each bar is one DFT bin, \(\SI{0.25}{\hertz}\) wide. \textbf{Band power} is
nothing more than the sum of the bins falling inside a band: the locomotor band
(\(0.5\text{--}3\,\)Hz, \(\approx 10\) bins) and the freeze band (\(3\text{--}8\,\)Hz,
\(\approx 20\) bins). The \(\SI{4}{\second}\) window is chosen precisely so these bins
are fine enough to separate the locomotor and freeze humps cleanly.}
\label{fig:bins}
\end{figure}

\begin{keyidea}
\textbf{Resolution costs time.} The bin spacing \(\Delta f = f_s/N\) shrinks as the
window grows. To resolve the \(3\text{--}8\,\)Hz freeze hump we need fine bins, so we
use a long \(N=256\) (\(\SI{4}{\second}\)) window: \(\SI{0.25}{\hertz}\) bins put
\(\sim\!20\) samples across the freeze band. A short window would give coarse, blurry
bins that smear walking and freezing together. The price is temporal: a \(4\,\)s window
cannot localise an event to better than a few seconds --- which is fine, because a
freeze \emph{lasts} seconds.
\end{keyidea}

\subsection{Power, leakage, and the Hann window}

We rarely care about the phase of each sinusoid, only how much \emph{energy} sits at
each frequency. That is the \textbf{power spectrum}, \(|X[k]|^2\): square the magnitude
of each coefficient. Squaring also makes the quantity additive in a way the next
section relies on. Parseval's theorem guarantees the bookkeeping is honest:
\begin{equation}
\sum_{n=0}^{N-1} a[n]^2 \;=\; \frac{1}{N}\sum_{k=0}^{N-1} \bigl|X[k]\bigr|^2 .
\label{eq:parseval}
\end{equation}

\begin{mathsbox}
Parseval says the total energy is the same whether you measure it in time
(left, sum of squared samples) or in frequency (right, sum of squared spectrum). No
energy is created or destroyed by changing viewpoint. This is why ``band power = sum
of \(|X[k]|^2\) over a band'' is a meaningful slice of the signal's real energy, and
why adding the locomotor and freeze band powers gives the leg's total movement energy
(used later by the still-floor gate).
\end{mathsbox}

There is a catch. The DFT secretly assumes the \(N\)-sample window repeats forever. A
real \(4\,\)s slice almost never starts and ends at the same value, so the assumed
loop has a discontinuity at the seam --- a sudden jump that the DFT explains by
sprinkling spurious energy across many bins. This smearing is called \textbf{spectral
leakage}, and it can wash freeze-band energy into the locomotor band (and vice versa),
exactly the contamination we must avoid.

The fix is a \textbf{window function} that gently tapers the slice to zero at both
ends before transforming, removing the seam. Cadence uses the \textbf{Hann window}:
\begin{equation}
w[n] \;=\; 0.5\Bigl(1-\cos\tfrac{2\pi n}{N-1}\Bigr),
\qquad n=0,1,\dots,N-1 .
\label{eq:hann}
\end{equation}
It rises smoothly from \(0\), peaks at the centre, and returns to \(0\) --- so however
the slice was cut, the windowed version always starts and ends flush.

\begin{figure}[H]
\centering
\resizebox{0.96\textwidth}{!}{%
\begin{tikzpicture}
% ----- LEFT: Hann window shape -----
\begin{axis}[
  name=hann,
  width=8.4cm,height=5.0cm,
  xlabel={sample index \(n\)}, ylabel={weight \(w[n]\)},
  axis lines=left, xmin=0,xmax=256, ymin=0,ymax=1.08,
  xtick={0,64,128,192,256}, ytick={0,0.5,1},
  tick label style={font=\footnotesize},
  every axis label/.append style={font=\small},
  title={\small\bfseries\color{ink} Hann window (\(N=256\))},
]
\addplot[teal,very thick,domain=0:256,samples=300]
  {0.5*(1-cos(deg(2*pi*x/255)))};
\addplot[draw=none,fill=paleteal,domain=0:256,samples=120]
  {0.5*(1-cos(deg(2*pi*x/255)))}\closedcycle;
\node[teal!60!black,font=\scriptsize,align=center] at (axis cs:128,0.45) {taper\\to zero\\at both\\ends};
\end{axis}
% ----- RIGHT: effect on the seam -----
\begin{axis}[
  at={($(hann.east)+(1.8cm,0)$)}, anchor=west,
  width=8.4cm,height=5.0cm,
  xlabel={sample index \(n\)}, ylabel={signal},
  axis lines=left, xmin=0,xmax=256, ymin=-1.25,ymax=1.25,
  xtick={0,128,256}, ytick=\empty,
  tick label style={font=\footnotesize},
  every axis label/.append style={font=\small},
  title={\small\bfseries\color{ink} Raw vs.\ windowed slice},
]
% raw signal: does not return to its start -> seam jump
\addplot[mute,thick,dashed,domain=0:256,samples=300]
  {sin(deg(2*pi*4.3*x/256))};
% windowed signal: tapered to zero at both ends
\addplot[navy,very thick,domain=0:256,samples=400]
  {0.5*(1-cos(deg(2*pi*x/255)))*sin(deg(2*pi*4.3*x/256))};
\node[mute,font=\scriptsize] at (axis cs:120,1.14) {raw (seam jump)};
\node[navy,font=\scriptsize,anchor=west] at (axis cs:14,-0.78) {Hann-tapered};
\end{axis}
\end{tikzpicture}}
\caption{\textbf{Left:} the Hann taper --- a raised cosine that weights the centre of
the window fully and fades the edges to zero. \textbf{Right:} multiplying the raw
slice (grey, dashed --- note it does \emph{not} return to its starting value, creating
a discontinuity when the DFT loops it) by the Hann window gives the tapered slice
(navy), which starts and ends flush at zero. The seam vanishes, and with it most of
the spectral leakage.}
\label{fig:hann}
\end{figure}

\subsection{Welch's method: a steadier spectrum}

A single \(4\,\)s DFT is noisy --- its bin heights jitter from one window to the next.
Cadence uses \textbf{Welch's method} to get a steadier estimate of the power spectral
density (PSD). The recipe is short:
\begin{enumerate}
\item Subtract the window mean from \(a[n]\) (kills the gravity/DC term so it does not
      dominate bin \(0\)).
\item Apply the Hann taper of Eq.~\eqref{eq:hann} (cuts leakage).
\item Take the FFT and form \(|X[k]|^2\) (the power at each bin).
\item Average these periodograms across overlapping windows (smooths the jitter).
\end{enumerate}
The result is the PSD: a stable estimate of power-per-frequency, normalised so that
summing it over a band yields that band's power. In code the whole chain is a few
lines built on the \code{fog.dsp} primitives:

\begin{lstlisting}[style=py,caption={The Welch PSD chain, then a band-power slice (from \code{fog.dsp}).}]
def welch_psd(window, fs=64, nfft=256):
    x = window - window.mean()          # 1. remove gravity / DC
    x = x * np.hanning(len(x))          # 2. Hann taper -> less leakage
    X = np.fft.rfft(x, n=nfft)          # 3. FFT (one-sided, real input)
    return (np.abs(X) ** 2) / fs        # power spectral density, |X|^2

def band_power(psd, fs, nfft, f_lo, f_hi):
    df  = fs / nfft                     # 0.25 Hz bin spacing
    lo  = int(np.ceil (f_lo / df))      # first bin in band
    hi  = int(np.floor(f_hi / df))      # last  bin in band
    return psd[lo:hi + 1].sum()         # band power = sum of bins
\end{lstlisting}

\begin{intuition}
Welch's method is the audio equivalent of letting a tuner ``settle'': instead of
trusting one fleeting snapshot, it averages several overlapping looks at the same few
seconds so random noise cancels out while the genuine walking and freeze tones stand
firm. The \(\SI{0.25}{\hertz}\) bins then read like a labelled ruler --- you can point
straight at the \(0.5\text{--}3\,\)Hz and \(3\text{--}8\,\)Hz stretches and add up
what is there.
\end{intuition}

Figure~\ref{fig:capture1-spectral} shows the payoff on real data. The
\code{capture1.csv} worksheet capture runs through four phases; in the freeze phase
the magnitude trace erupts into a fast, dense tremble, and the Freeze Index --- a
ratio of two band powers we have just learned to compute --- jumps from its calm
baseline of \(\approx 1\) to a peak of \(21.35\).

\begin{figure}[H]
\centering
\includegraphics[width=0.82\linewidth]{capture1_trace}
\caption{Real worksheet capture (\code{capture1.csv}, \SI{64}{\hertz},
\SI{41.4}{\second}). \textbf{Top:} the ankle magnitude \(\norm{a}\) in mg across four
phases. The calm phases are a quiet wobble near \(1\,g\); the freeze phase (phase~2)
explodes into a fast, high-amplitude tremble. \textbf{Bottom:} the Freeze Index
computed per window from the band powers defined here --- still \(\approx 1.1\), walk
\(\approx 1.2\), then a freeze peak of \(21.35\). The fast tremble is invisible as
``size'' alone but obvious once we read it in the frequency domain.}
\label{fig:capture1-spectral}
\end{figure}

\begin{takeaway}
Stop measuring \emph{how big} the ankle wobbles and measure \emph{how fast}: the DFT
turns a \(\SI{4}{\second}\) window into \(\SI{0.25}{\hertz}\) bins, the Hann-tapered
Welch PSD makes those bins trustworthy, and walking (\(0.5\text{--}3\,\)Hz) lands
elsewhere than freezing (\(3\text{--}8\,\)Hz) --- the foundation of the Freeze Index.
\end{takeaway}
